% ============================================================================
% Paper template entry point
% Compile with: latexmk -xelatex -interaction=nonstopmode -halt-on-error main.tex
% ============================================================================
\documentclass[11pt,a4paper]{article}

% Shared packages, layout, theorem environments, and reusable macros.
% The template is intentionally venue-neutral. Replace the document class and
% citation style in main.tex when a target venue provides its own style file.

% ---- Encoding and typography ------------------------------------------------
% XeLaTeX/LuaLaTeX uses ctex so that the Chinese drafting guidance below is
% rendered correctly. The pdfLaTeX branch keeps the template usable for an
% English-only manuscript.
\usepackage{iftex}
\ifPDFTeX
  \usepackage[T1]{fontenc}
  \usepackage[utf8]{inputenc}
  \usepackage{lmodern}
\else
  % Fandol is shipped with TeX Live and avoids relying on system-specific
  % CJK fonts (for example STHeiti on macOS).
  \usepackage[fontset=fandol]{ctex}
\fi
\usepackage{microtype}
\usepackage{xspace}

% ---- Page layout ------------------------------------------------------------
\usepackage[a4paper,margin=1in]{geometry}
\usepackage{setspace}
\setstretch{1.08}
\usepackage{parskip}
\setlength{\parindent}{1.5em}
\usepackage{placeins}

% ---- Mathematics, algorithms, and units -----------------------------------
\usepackage{amsmath,amssymb,amsthm}
\usepackage{mathtools}
\usepackage{siunitx}
\usepackage{algorithm}
\usepackage{algpseudocode}

% ---- Figures, tables, and lists --------------------------------------------
\usepackage{graphicx}
\graphicspath{{figures/}{figures/main/}{figures/appendix/}}
\usepackage{booktabs}
\usepackage{multirow}
\usepackage{tabularx}
\usepackage{array}
\usepackage{makecell}
\usepackage{subcaption}
\usepackage{float}
\usepackage{enumitem}
\usepackage{xcolor}

% ---- Citations and links ----------------------------------------------------
\usepackage[round,authoryear]{natbib}
\usepackage{xurl}
\usepackage{hyperref}
\usepackage[nameinlink,noabbrev]{cleveref}

\hypersetup{
  colorlinks=true,
  linkcolor=blue!55!black,
  citecolor=green!40!black,
  urlcolor=blue!65!black,
  pdftitle={Your Paper Title}
}

% ---- Colors and writing helpers --------------------------------------------
\definecolor{todoRed}{HTML}{B00020}
\definecolor{noteBlue}{HTML}{0B5394}
\definecolor{lightGray}{HTML}{F3F4F6}

% Visible markers make unfinished sections easy to find during drafting.
% Remove or redefine this command before submission if desired.
\newcommand{\placeholder}[1]{\textcolor{todoRed}{\textbf{[TODO: #1]}}}
\newcommand{\note}[1]{\textcolor{noteBlue}{\textit{[Note: #1]}}}
\newcommand{\methodname}{\textsc{MethodName}\xspace}
\newcommand{\datasetname}{\textsc{DatasetName}\xspace}
\newcommand{\eg}{e.g.,\xspace}
\newcommand{\ie}{i.e.,\xspace}

% A blank figure placeholder that compiles before real assets are available.
\newcommand{\placeholderfigure}[2]{%
  \begin{figure}[t]
    \centering
    \fcolorbox{black!35}{lightGray}{%
      \rule{0pt}{1.35in}\rule{0.88\linewidth}{0pt}%
    }
    \caption{#1}
    \label{#2}
  \end{figure}%
}

% ---- Theorem-like environments --------------------------------------------
\theoremstyle{plain}
\newtheorem{theorem}{Theorem}
\newtheorem{proposition}{Proposition}
\newtheorem{lemma}{Lemma}
\newtheorem{corollary}{Corollary}
\theoremstyle{definition}
\newtheorem{definition}{Definition}
\newtheorem{example}{Example}
\theoremstyle{remark}
\newtheorem{remark}{Remark}

% ---- Table helpers ----------------------------------------------------------
\newcolumntype{Y}{>{\centering\arraybackslash}X}
\newcommand{\best}[1]{\textbf{#1}}
\newcommand{\secondbest}[1]{\underline{#1}}

% ---- Consistent figure/table references ------------------------------------
\crefname{algorithm}{Algorithm}{Algorithms}
\crefname{figure}{Figure}{Figures}
\crefname{table}{Table}{Tables}
\crefname{equation}{Equation}{Equations}


% ---- Metadata: edit these fields first -------------------------------------
\newcommand{\papertitle}{Your Paper Title: A Clear and Specific Subtitle}
\newcommand{\paperauthor}{First Author \\ Institution \\ \texttt{first.author@example.com}}
\newcommand{\paperkeywords}{keyword one; keyword two; keyword three}

\title{\papertitle}
\author{\paperauthor}
\date{\today}

\begin{document}

\maketitle

\begin{abstract}
% 用 150--250 字回答四个问题：研究什么、为什么重要、怎么做、得到什么结果。
\placeholder{写摘要：问题背景、方法、核心结果和结论。}
\end{abstract}

\paragraph{Keywords:} \paperkeywords

% ---- Main text --------------------------------------------------------------
\section{Introduction}
\label{sec:introduction}

% Recommended flow:
% (1) establish the broad problem and its importance;
% (2) identify the concrete gap or failure mode;
% (3) state the proposed idea at a high level;
% (4) preview the evidence;
% (5) list contributions that can be checked in the paper.

\paragraph{Problem and motivation.}
\placeholder{用一段话说明研究问题、应用背景和实际影响。}

\paragraph{Research gap.}
\placeholder{说明现有方法在哪个具体场景下不足，以及为什么现有工作没有解决它。}

\paragraph{Approach and preview.}
\placeholder{用非技术语言概述你的方法，并给出最重要的实验发现。}

\paragraph{Contributions.}
Our work makes the following contributions:
\begin{itemize}[leftmargin=2em]
  \item \placeholder{贡献一：提出什么问题、数据集、理论或系统。}
  \item \placeholder{贡献二：方法中最关键的技术创新。}
  \item \placeholder{贡献三：用哪些实验或分析验证了什么结论。}
\end{itemize}

\note{每条贡献都应在正文中有对应的章节、表格、图或定理支撑。}

\section{Related Work}
\label{sec:related-work}

% Organize by research themes rather than by a paper-by-paper chronology.
% For each theme: summarize the consensus, identify the limitation, and explain
% how the present work differs. Replace the sample citation with your sources.

\subsection{Research area one}
\placeholder{概括第一个相关方向的代表性方法、发展脉络和不足。}

\subsection{Research area two}
\placeholder{概括第二个相关方向，并明确与本文的关系。}

\subsection{Positioning of this work}
\placeholder{用一段话总结本文与最接近工作的差异，避免只罗列引用。}

\note{引用示例：叙述式引用使用
\texttt{\textbackslash citet\{key\}}，括号式引用使用
\texttt{\textbackslash citep\{key\}}。}

\section{Method}
\label{sec:method}

% Start with an overview, then define notation, describe each component, and
% finish with complexity, implementation choices, or theoretical guarantees.

\subsection{Overview}
\placeholder{给出方法总览。建议先放一张系统流程图，再说明输入、核心模块和输出。}

% Replace the following placeholder with an actual asset when available:
% \includegraphics[width=0.92\\linewidth]{figures/method_overview.pdf}
% \caption{Overview of the proposed method.}
% \label{fig:method-overview}
\placeholderfigure{Method overview placeholder. Replace the box with a figure in the \texttt{figures/} directory.}{fig:method-overview}

\subsection{Problem formulation}
Let $\mathcal{X}$ denote the input space, $\mathcal{Y}$ the output space, and
$f_\theta: \mathcal{X}\to\mathcal{Y}$ the proposed model or procedure. Given a
training set $\mathcal{D}=\{(x_i,y_i)\}_{i=1}^{n}$, the generic objective is
\begin{equation}
  \min_{\theta}\; \frac{1}{n}\sum_{i=1}^{n}
  \mathcal{L}\bigl(f_\theta(x_i),y_i\bigr)+\lambda\,\Omega(\theta),
  \label{eq:objective}
\end{equation}
where $\mathcal{L}$ is the task loss, $\Omega$ is an optional regularizer, and
$\lambda$ controls its strength. Replace this generic formulation with the
notation and objective specific to your work.

\subsection{Core component}
\placeholder{定义最关键的模块、公式、假设和设计选择。解释每个符号及其作用。}

\begin{algorithm}[t]
  \caption{Generic training or inference procedure}
  \label{alg:generic-procedure}
  \begin{algorithmic}[1]
    \Require Dataset $\mathcal{D}$, initialization $\theta_0$, number of steps $T$
    \Ensure Trained parameters or predictions
    \State $\theta \gets \theta_0$
    \For{$t=1,\ldots,T$}
      \State Sample a mini-batch $B\subseteq\mathcal{D}$
      \State Compute objective $\mathcal{L}_B$ using \cref{eq:objective}
      \State Update $\theta$ using the chosen optimizer or rule
    \EndFor
    \State \Return $\theta$
  \end{algorithmic}
\end{algorithm}

\subsection{Analysis and implementation details}
\placeholder{说明复杂度、理论性质、超参数、工程实现和可复现性细节。}

\section{Experimental Setup}
\label{sec:experiments}

\subsection{Research questions}
State the questions before presenting results. For example:
\begin{description}[leftmargin=2.4em,style=nextline]
  \item[RQ1:] \placeholder{方法是否优于强基线？}
  \item[RQ2:] \placeholder{每个模块分别带来多少收益？}
  \item[RQ3:] \placeholder{方法在哪些数据、规模或条件下失效？}
\end{description}

\subsection{Datasets and preprocessing}
\placeholder{介绍数据来源、规模、划分、许可、预处理和数据泄漏控制。}

\subsection{Baselines and evaluation metrics}
\placeholder{说明基线选择、公平比较方式、评价指标定义及统计检验。}

\subsection{Implementation details}
\placeholder{报告硬件、软件版本、训练轮数、优化器、学习率、随机种子和调参范围。}

\begin{table}[t]
  \centering
  \caption{Template for dataset or experimental configuration statistics.}
  \label{tab:dataset-statistics}
  \begin{tabular}{lrrr}
    \toprule
    Split & Instances & Features & Labels \\
    \midrule
    Train & \placeholder{--} & \placeholder{--} & \placeholder{--} \\
    Validation & \placeholder{--} & \placeholder{--} & \placeholder{--} \\
    Test & \placeholder{--} & \placeholder{--} & \placeholder{--} \\
    \bottomrule
  \end{tabular}
\end{table}

\section{Results and Discussion}
\label{sec:results}

% Present one claim per subsection. Start each subsection with the conclusion,
% then provide the table/figure evidence, followed by interpretation.

\subsection{Main results}
\placeholder{先用一句话给出结论，再引用主结果表并解释差异。}

\begin{table}[t]
  \centering
  \caption{Template for the main comparison. Bold and underlined values indicate the best and second-best results.}
  \label{tab:main-results}
  \begin{tabular}{lccc}
    \toprule
    Method & Metric 1 $\uparrow$ & Metric 2 $\uparrow$ & Runtime $\downarrow$ \\
    \midrule
    Baseline A & \placeholder{--} & \placeholder{--} & \placeholder{--} \\
    Baseline B & \placeholder{--} & \placeholder{--} & \placeholder{--} \\
    \methodname & \best{\placeholder{--}} & \secondbest{\placeholder{--}} & \placeholder{--} \\
    \bottomrule
  \end{tabular}
\end{table}

\subsection{Ablation and sensitivity analysis}
\placeholder{逐项移除模块或改变关键设置，说明性能变化是否支持方法设计。}

\subsection{Qualitative or error analysis}
\placeholder{给出典型成功、失败案例，并总结可解释的错误模式。}

\subsection{Limitations and threats to validity}
\placeholder{讨论外部有效性、数据偏差、计算成本、测量误差和可推广边界。}

\section{Conclusion}
\label{sec:conclusion}

\placeholder{用一到两段总结问题、方法、核心发现和实际意义。}

\paragraph{Future work.}
\placeholder{提出由本文结果直接支持的后续研究方向，避免泛泛而谈。}


% ---- References -------------------------------------------------------------
% 使用 natbib 命令：\citet{key}（叙述式）和 \citep{key}（括号式）。
% \nocite{*} 让空白模板也能生成示例参考文献；开始正式写作后可删除它。
\nocite{*}
\bibliographystyle{plainnat}
\bibliography{references}

% ---- Supplementary material -------------------------------------------------
\clearpage
\appendix
\section{Additional Method Details}
\label{app:method-details}

\placeholder{放置正文中会打断叙事，但对复现、证明或审稿很重要的细节。}

\subsection{Proofs or derivations}
\placeholder{补充定理证明、推导过程或复杂度分析。}

\begin{lemma}
  Replace this statement with a lemma used by the main text.
\end{lemma}

\begin{proof}
  \placeholder{补充证明。}
\end{proof}

\section{Additional Experimental Results}
\label{app:additional-results}

\placeholder{放置额外数据集、完整结果、提示词、超参数、用户研究问卷或案例。}

\begin{table}[H]
  \centering
  \caption{Template for a supplementary result table.}
  \label{tab:supplementary-results}
  \begin{tabular}{lcc}
    \toprule
    Setting & Mean $\pm$ Std. & Number of runs \\
    \midrule
    Setting A & \placeholder{--} & \placeholder{--} \\
    Setting B & \placeholder{--} & \placeholder{--} \\
    \bottomrule
  \end{tabular}
\end{table}


\end{document}
